The organization \must{} maintain an accurate inventory of technology assets that are owned, leased, managed, operated, or used to access organizational systems or data.

Each asset \must{} have a unique identifier that allows it to be distinguished from other assets throughout its lifecycle. Physical assets \should{} be labeled with an organizational asset tag where practical.

The asset inventory \must{} include, where applicable:
\begin{itemize}
    \item Servers, workstations, laptops, tablets, and mobile devices
    \item Virtual machines, containers, hypervisors, and cloud resources
    \item Firewalls, routers, switches, wireless access points, and VPN equipment
    \item Storage systems, backup systems, and removable media
    \item Printers, scanners, cameras, and other network connected devices
    \item Internet facing systems and externally hosted services
    \item Devices used by employees, contractors, vendors, and other authorized users
    \item Equipment used to store, process, transmit, or access organizational data
\end{itemize}

Each inventoried asset \must{}, to the extent technically and operationally feasible, have the following information recorded:
\begin{itemize}
    \item Asset name or hostname
    \item Asset type and description
    \item Hardware manufacturer, model, and serial number
    \item Assigned user, department, or responsible role
    \item Business or system owner
    \item Physical location or hosting environment
    \item Internet Protocol address, media access control address, or other identifier
    \item Operating system or firmware, where applicable
    \item Data classification and types of data handled
    \item Criticality and recovery priority
    \item Related system, service, or business process
    \item Support status and warranty information
    \item Acquisition date and expected replacement date
    \item Current status, such as active, inactive, lost, under repair, or retired
    \item Date of last inventory review
\end{itemize}

Asset owners \must{} ensure that asset information remains accurate and current. New assets \must{} be recorded before being placed into production or connected to organizational networks, unless an emergency deployment is authorized and the asset is documented as soon as practicable.

Assets \must{} be assigned to a responsible owner or role. The owner \must{} be accountable for the asset's approved use, security requirements, maintenance, and retirement.

Assets \must{} be classified according to their importance and the sensitivity of the data they store, process, or transmit. Critical assets \must{} be identified so that they can receive appropriate protection, monitoring, backup, and recovery planning.

Unauthorized, unknown, or unmanaged assets \must{} be investigated. An asset that cannot be approved or associated with a responsible owner \must{} be isolated, removed, or otherwise restricted until its status is resolved.

The organization \should{} use automated discovery, endpoint management, cloud management, network monitoring, or similar technical controls to identify assets that are missing from the inventory.

The asset inventory \must{} be reviewed:
\begin{itemize}
    \item At least annually
    \item After significant changes to systems or infrastructure
    \item During major hardware or software deployments
    \item After mergers, acquisitions, office moves, or organizational changes
    \item Following the loss, theft, or disposal of equipment
    \item Following a security incident that may affect asset ownership or status
\end{itemize}

Critical assets \should{} be reviewed more frequently based on their risk and operational importance.

Asset records \must{} be updated when equipment is reassigned, relocated, modified, replaced, lost, stolen, or retired. Changes to ownership, responsible personnel, network location, operating system, or data classification \must{} be recorded promptly.

Lost or stolen assets \must{} be reported to management or the designated security contact promptly. The organization \must{}, where technically possible, disable accounts, revoke credentials, remove access tokens, enable remote lock or wipe, and assess whether organizational data may have been exposed.

Assets that are no longer required \must{} be securely decommissioned. Before disposal, reuse, return, or transfer, organizational data \must{} be removed using an approved sanitization method appropriate to the media and sensitivity of the information.

Asset disposal or retirement records \should{} include:
\begin{itemize}
    \item Asset identifier
    \item Date of disposal or retirement
    \item Reason for disposal or retirement
    \item Sanitization or destruction method
    \item Person or vendor performing the work
    \item Transfer, recycling, return, or destruction details
    \item Approval, where required
\end{itemize}

Third-party or cloud-hosted assets that store, process, or transmit organizational data \must{} be included in the inventory or linked to an inventory record. The organization \must{} maintain sufficient information to identify the service provider, service owner, data handled, access method, and contract or subscription status.

The asset inventory \should{} be protected from unauthorized modification or disclosure. Access \must{} be limited to personnel with a legitimate business need, and changes to critical inventory records \should{} be logged.